% !TEX root = master.tex
% !TEX program = lualatex
\lecture{6}{26-03-2026}{Residues and Residue Theorem}{}
\section{Residues}

\begin{definition}[Residues]
 
Let $f$ be holomorphic on a punctured neighborhood of $z_0$, and write its Laurent series at $z_0$ as
\[
f(z)=\sum_{n=-\infty}^{\infty} c_n (z-z_0)^n.
\]
The coefficient $c_{-1}$ of the term $(z-z_0)^{-1}$ is called the \emph{residue} of $f$ at $z_0$:
\[
\operatorname{Res}_{z_0}(f)=c_{-1}
\]

\end{definition}
\subsection{Integral expression of the residue}

By the definition of the coefficients in the Laurent series,
\[
\operatorname{Res}_{z_0}(f)
=
\frac{1}{2\pi i}\int_{\gamma} f(z)\,dz,
\]
where $\gamma$ is a closed, simple, counterclockwise oriented curve that contains $z_0$ in its interior, and such that $f$ is holomorphic on $int_{\gamma}\setminus \{z_0\}$.

\paragraph{Example}

If $f$ is holomorphic on an open set containing $z_0$, then
\[
\operatorname{Res}_{z_0}(f)=0
\]
by Cauchy's theorem.

For
\[
f(z)=\frac{1}{z},
\]
since it's already the Laurent series at $0$, hence
\[
\operatorname{Res}_0\!\left(\frac{1}{z}\right)=1
\]

\subsection{Formula for a pole of order \texorpdfstring{$m$}{m}}

Suppose that $f$ has a pole of order $m$ at $z_0$, so that its Laurent series is
\[
f(z)=c_{-m}(z-z_0)^{-m}+\cdots+c_{-1}(z-z_0)^{-1}+c_0+c_1(z-z_0)+\cdots
\]
Consider the holomorphic function
\[
g(z)=(z-z_0)^m f(z) =c_{-m}+c_{-m+1}(z-z_0)+\cdots+c_{-1}(z-z_0)^{m-1}+c_0(z-z_0)^m+\cdots
\]
Therefore,
\begin{corollary}
  
\[
c_{-1}= \lim_{z \to z_0} \frac{1}{(m-1)!}\frac{d^{\,m-1}}{dz^{m-1}}\Big((z-z_0)^m f(z))\Big)
\]

This formula is useful when we already know that $f$ has a pole of order $m$ at $z_0$.

In particular, if $m=1$, then
\[
\operatorname{Res}_{z_0}(f)
=
\lim_{z\to z_0}(z-z_0)f(z).
\]

\end{corollary}

\paragraph{Example}

\[
f(z)=\frac{1}{\sin z}
\]
At $z_0=0$, this function has a pole of order $1$. 

Hence,
\[
\operatorname{Res}_0(f)
=
\lim_{z\to 0} z\,f(z)
=
\lim_{z\to 0}\frac{z}{\sin z}
\]
Using
\[
\sin z = z - \frac{z^3}{3!}+\cdots,
\]
we get
\[
\frac{z}{\sin z}
=
\frac{z}{z-\frac{z^3}{3!}+\cdots}
=
\frac{1}{1-\frac{z^2}{3!}+\cdots},
\]
so the limit is
\[
\operatorname{Res}_0(f)=1.
\]

\section{Residue theorem}

\begin{theorem}[Residue theorem]
Suppose $D\subseteq \mathbb{C}$ is open and simply connected, and let $\gamma$ be a simple closed curve in $D$ which is counterclockwise oriented. Let $z_1,\dots,z_m\in \operatorname{int}(\gamma)$. Assume
\[
f: D\setminus\{z_1,\dots,z_m\}\to \mathbb{C}
\]
is holomorphic. Then
\[
\int_\gamma f(z)\,dz
=
2\pi i \sum_{i=1}^{m}Res_{z_i}(f) 
\]
\end{theorem}

Hence, in order to compute $\int_\gamma f(z)\,dz$, it suffices to compute the residues of the singular points in the interior of $\gamma$.

\subsection{Idea of the proof}

For simplicity, assume there are only two singular points $z_1$ and $z_2$. One deforms the contour $\gamma$ into two small positively oriented curves $\gamma_1$ and $\gamma_2$ around the singularities, connected by a very narrow bridge. By Cauchy's theorem, the integral over $\gamma$ is equal to the sum of the integrals over the small circles, since the contributions of the bridge cancel. Therefore,
\[
\int_\gamma f(z)\,dz
=
\int_{\gamma_1} f(z)\,dz
+
\int_{\gamma_2} f(z)\,dz
=
2\pi i\,\operatorname{Res}_{z_1}(f)
+
2\pi i\,\operatorname{Res}_{z_2}(f).
\]

\paragraph{Example}

\[
f(z)=\frac{1}{z}.
\]
Let $\gamma$ be a simple closed counterclockwise oriented curve.

\begin{enumerate}
    \item If $0\in \operatorname{int}(\gamma)$, then by the residue theorem,
    \[
    \int_\gamma \frac{1}{z}\,dz = 2\pi i\,\operatorname{Res}_0\!\left(\frac{1}{z}\right)=2\pi i.
    \]

    \item If $0\notin \operatorname{int}(\gamma)$, then
    \[
    \int_\gamma \frac{1}{z}\,dz=0.
    \]

    \item If $0\in \gamma$, then the integral is not well-defined.
\end{enumerate}

\paragraph{Example}
\[
f(z)=\frac{1}{z(z+1)}.
\]
We have two simple poles, at $z=0$ and $z=-1$.

Their residues are
\[
\operatorname{Res}_0(f)
=
\lim_{z\to 0} zf(z)
=
\lim_{z\to 0}\frac{1}{z+1}
=1,
\]
and
\[
\operatorname{Res}_{-1}(f)
=
\lim_{z\to -1}(z+1)f(z)
=
\lim_{z\to -1}\frac{1}{z}
=-1.
\]

So:

\begin{enumerate}
    \item If neither $0$ nor $-1$ lies in $\operatorname{int}(\gamma)$, then
    \[
    \int_\gamma f(z)\,dz = 0.
    \]

    \item If $0\in \operatorname{int}(\gamma)$ but $-1\notin \operatorname{int}(\gamma)$, then
    \[
    \int_\gamma f(z)\,dz = 2\pi i.
    \]

    \item If $-1\in \operatorname{int}(\gamma)$ but $0\notin \operatorname{int}(\gamma)$, then
    \[
    \int_\gamma f(z)\,dz = -2\pi i.
    \]

    \item If both $0$ and $-1$ lie in $\operatorname{int}(\gamma)$, then
    \[
    \int_\gamma f(z)\,dz
    =
    2\pi i\big(\operatorname{Res}_0(f)+\operatorname{Res}_{-1}(f)\big)
    =
    2\pi i(1-1)=0.
    \]

    \item If $0\in\gamma$ or $-1\in\gamma$, then the integral is not well-defined.
\end{enumerate}

\paragraph{Example}

$f: \mathbb{C}\setminus \{0,1\} \to \mathbb{C}$,

\[
f(z)=\frac{1}{z^2}+\frac{2}{z}+\frac{3}{z-1}.
\]
Then $f$ has a pole of order $2$ at $0$ and a pole of order $1$ at $1$.

At $z=1$, since the pole is simple,
\[
\operatorname{Res}_1(f)
=
\lim_{z\to 1}(z-1)f(z)
=
\lim_{z\to 1}\left(\frac{z-1}{z^2}+\frac{2(z-1)}{z}+3\right)
=3.
\]

At $z=0$, we use the formula for a pole of order $2$:
\[
\operatorname{Res}_0(f)
= \lim_{z \to 0} \frac{1}{1!}
\frac{d}{dz}\big(z^2 f(z)\big).
\]
Now
\[
z^2 f(z)
=
1+2z+\frac{3z^2}{z-1},
\]
so
\[
\frac{d}{dz}\big(z^2 f(z)\big)
=
2+\frac{6z}{z-1}-\frac{3z^2}{(z-1)^2}.
\]
Evaluating at $z=0$ gives
\[
\operatorname{Res}_0(f)=2.
\]

Actually this could have been seen by reading directly $f(z)$.

If $\gamma(t)=10e^{it}$ for $t\in[0,2\pi]$, then both poles lie inside $\gamma$, hence
\[
\int_\gamma f(z)\,dz
=
2\pi i\big(\operatorname{Res}_0(f)+\operatorname{Res}_1(f)\big)
=
2\pi i(2+3)
=
10\pi i.
\]

\subsection{Useful remark}

It is often advisable to first rewrite the function in a form where the residue can be read off immediately.

For example:
\[
e^{1/z}=1+\frac{1}{z}+\frac{1}{2!}\frac{1}{z^2}+\cdots
\quad\Longrightarrow\quad
\operatorname{Res}_0\!\left(e^{1/z}\right)=1.
\]

On the other hand,
\[
e^{1/z^2}=1+\frac{1}{z^2}+\frac{1}{2!}\frac{1}{z^4}+\cdots
\quad\Longrightarrow\quad
\operatorname{Res}_0\!\left(e^{1/z^2}\right)=0.
\]
